\begin{frame}{模板说明}
This template is a secondary creation of \bhref{https://www.overleaf.com/latex/templates/sintef-presentation/jhbhdffczpnx}{SINTEF Presentation} template from \bhref{mailto:federico.zenith@sintef.no}{Federico Zenith}. All rights reserved by him.

This template is released under the \bhref{https://creativecommons.org/licenses/by-nc/4.0/legalcode}{Creative Commons CC BY 4.0} license.
\end{frame}

\section{Introduction}

\begin{frame}{Beamer for SINTEF slides}
\begin{itemize}
\item We assume you can use \LaTeX; if you cannot, \bhref{http://en.wikibooks.org/wiki/LaTeX/}{you can learn it here}.
\item Beamer is one of the most popular and powerful document classes for presentations in \LaTeX.
\item Beamer has a detailed \bhref{http://www.ctan.org/tex-archive/macros/latex/contrib/beamer/doc/beameruserguide.pdf}{user manual}.
\item Here we present only the most basic features to get you up to speed.
\end{itemize}
\end{frame}

\begin{frame}{Beamer vs. PowerPoint}
Compared to PowerPoint, using \LaTeX\ is better because:
\begin{itemize}
\item It is not What-You-See-Is-What-You-Get, but What-You-\emph{Mean}-Is-What-You-Get: you write the content, the computer does the typesetting.
\item It produces a \texttt{pdf}: no problems with fonts, formulas, or program versions.
\item It is easier to keep consistent style, fonts, and highlighting.
\item Math typesetting in \TeX\ is the best:
\begin{equation*}
\mathrm{i}\,\hslash\frac{\partial}{\partial t}\Psi(\mathbf{r},t)=-\frac{\hslash^2}{2m}\nabla^2\Psi(\mathbf{r},t)+V(\mathbf{r})\Psi(\mathbf{r},t)
\end{equation*}
\end{itemize}
\end{frame}

\section{Editing}

\begin{frame}[fragile]{Selecting the Class}
After the last update to the graphic profile, the \texttt{sintef} theme for Beamer has been updated into a full-fledged class.

To start working with \texttt{sintefbeamer}, start a \LaTeX\ document with the preamble:
\begin{block}{Minimum Beamer Document}
\begin{lstlisting}[language=TeX]
\documentclass{beamer}
\begin{document}

\begin{frame}{Hello, world!}
% contents
\end {frame}

\end {document}
\end{lstlisting}
\end{block}
\end{frame}

\begin{frame}[fragile]{Title page}
To set a typical title page, you call some commands in the preamble:
\begin{block}{The Commands for the Title Page}
\begin{lstlisting}[language=TeX]
\title{Sample Title}
\subtitle{Sample subtitle}
\author{First Author, Second Author}
\date{Defaults to today's}
\end{lstlisting}
\end{block}
You can then write out the title page with \verb|\maketitle|.

You can set a different background image than the default one with the \verb|\titlebackground| command, set before \verb|\maketitle|.
\end{frame}

\begin{frame}[fragile]{Writing a Simple Slide}
\framesubtitle{It's really easy!}
\begin{itemize}[<+->]
\item A typical slide has bulleted lists.
\item These can be uncovered in sequence.
\end{itemize}
\begin{block}{Code for a Page with an Itemised List}<+->
\begin{lstlisting}[language=TeX]
\begin{frame}{Writing a Simple Slide}
\framesubtitle{It's really easy!}
\begin{itemize}[<+->]
\item A typical slide has bulleted lists.
\item These can be uncovered in sequence.
\end{itemize}
\end {frame}
\end{lstlisting}
\end{block}
\end{frame}

\begin{frame}[fragile]{Adding images}
Adding images works like in normal \LaTeX:
\begin{columns}
\begin{column}{0.7\textwidth}
\begin{block}{Code for Adding Images}
\begin{lstlisting}[language=TeX]
\includegraphics[width=\textwidth]{figs/img.jpg}
\end{lstlisting}
\end{block}
\end{column}
\begin{column}{0.3\textwidth}
\includegraphics[width=\textwidth]{figs/img.jpg}
\end{column}
\end{columns}
\end{frame}

\begin{frame}[fragile]{Splitting in Columns}
Splitting the page is easy and common; typically, one side has a picture and the other text:
\begin{columns}
\begin{column}{0.6\textwidth}
This is the first column.
\end{column}
\begin{column}{0.3\textwidth}
This is the second column.
\end{column}
\end{columns}
\begin{block}{Column Code}
\begin{lstlisting}[language=TeX]
\begin{columns}
\begin{column}{0.6\textwidth}
This is the first column.
\end{column}
\begin{column}{0.3\textwidth}
This is the second column.
\end{column}
\end {columns}
\end{lstlisting}
\end{block}
\end{frame}

\begin{frame}{Fonts}
\begin{itemize}
\item The paramount task of fonts is being readable.
\item There are good ones:
\begin{itemize}
\item {\textrm{Use serif fonts only with high-definition projectors.}}
\item {\textsf{Use sans-serif fonts otherwise, or if you simply prefer them.}}
\end{itemize}
\item There are not so good ones:
\begin{itemize}
\item {\texttt{Never use monospace for normal text.}}
\item {\frakfamily Gothic, calligraphic, or weird fonts should always be avoided.}
\end{itemize}
\end{itemize}
\end{frame}

\begin{frame}[fragile]{Look}
\begin{itemize}
\item To change between the light and dark themes, give the class options \texttt{light} or \texttt{dark}.
\item To insert a final slide, use \verb|\QApage|.
\item The aspect ratio is fixed at 16:9 by the theme.
\end{itemize}
\end{frame}

\section{Summary}

\begin{frame}{Good Luck!}
\begin{itemize}
\item Enough for an introduction. You should know enough by now.
\item If you have corrections or suggestions, \bhref{mailto:federico.zenith@sintef.no}{send them to me}.
\end{itemize}
\end{frame}

\begin{frame}{Block examples}
\begin{block}{block}
block
\end{block}
\begin{exampleblock}{example}
example
\end{exampleblock}
\begin{alertblock}{alert}
alert
\end{alertblock}
\end{frame}

